% Generated draft: requires mathematical and source review.
\documentclass[11pt]{article}
\usepackage[margin=25mm]{geometry}
\usepackage{fontspec}
\setmainfont{Latin Modern Roman}
\usepackage{amsmath,amssymb,mathtools,unicode-math}
\setmathfont{Latin Modern Math}
\usepackage{enumitem,needspace,xurl,hyperref}
\setlength{\parindent}{0pt}
\setlength{\parskip}{5pt}
\setlength{\emergencystretch}{3em}
\newcommand{\N}{\mathbb{N}}
\newcommand{\Z}{\mathbb{Z}}
\newcommand{\Q}{\mathbb{Q}}
\newcommand{\R}{\mathbb{R}}
\newcommand{\C}{\mathbb{C}}
\newcommand{\F}{\mathbb{F}}
\newcommand{\A}{\mathbb{A}}
\newcommand{\PP}{\mathbb{P}}
\newcommand{\E}{\mathbb{E}}
\newcommand{\eps}{\varepsilon}
\newcommand{\dd}{\,\mathrm{d}}
\begin{document}
\section*{UnsolvedMath ID: 4000017}
\textbf{Source problem code:} catalogue entry
\subsection{Definitions and mathematical statement}
Let $(X,d)$ be a metric space with a probability transition law $m_x$ at each point. For $x\ne y$, Ollivier coarse Ricci curvature is
\[
\kappa(x,y)=1-\frac{W_1(m_x,m_y)}{d(x,y)},
\]
where $W_1$ is the minimum expected distance over probability couplings, assuming finite first moments. The proposed scalar-curvature candidate is
\[
S(x)=\int_X\kappa(x,y)\,\mathrm{d}m_x(y),
\]
possibly replaced by a distance-weighted average
\[
S_w(x)=\int_X w(d(x,y))\kappa(x,y)\,\mathrm{d}m_x(y).
\]
These formulas require integrability. The unweighted expression also requires a convention at $y=x$ if $m_x$ has an atom there, because the displayed coarse-curvature formula is undefined on the diagonal. The source gives no diagonal convention or particular weight $w$.

The problem asks whether this averaged quantity has useful geometric or probabilistic properties, with volume-growth control offered as an example. A volume-growth question would concern the mass of metric balls $B(x,r)=\{y:d(x,y)\leq r\}$ with respect to a specified global measure, such as counting measure on a locally finite graph. The source supplies neither that global measure nor hypotheses and a quantitative volume comparison. Accordingly this is a proposal to develop and test a scalar-curvature notion, not a stated universal volume bound.
\subsection{Short English statement}
Study the average of coarse Ricci curvature over a point's transition law, possibly weighted by distance. Does this proposed scalar curvature control useful geometry or probability, for example growth of metric balls?
\subsection{Sources}
\begin{itemize}
\item Local UnsolvedMath registry, permanent ID 4000017, exact source code catalogue entry, ``Discrete scalar curvature''; proposed by Yann Ollivier, 2008. Registry status open, classification OPEN-TRIAGE.
\item Yann Ollivier, \emph{Discrete Ricci curvature: Open problems} (2008), Problem Q, \url{https://www.yann-ollivier.org/rech/publs/problems_curvmarkov.pdf}. Supplied pdf-text extraction accessed July 29, 2026; the item is presented as open, and no later authoritative resolution check is recorded. No PDF was accessed here.
\item The supplied research summary describes the functional as a proposal and reports no verified established theory or volume-growth control for this particular quantity. This is a limited supplied search assessment, not a claim that no such result exists.
\item Provenance flags status and rights NEEDS\_REVIEW: public access was verified, but redistribution permission was not established.
\end{itemize}
\end{document}
